\section{Predicate Schema and Strict Verifier}
\label{sec:predicate}

The bilateral admission envelope is a DSSE statement~\cite{dsse} carrying the predicate type \codepath{chio.bilateral-cosign-invocation.v1}. The payload type is the in-toto statement type, the statement type is in-toto v1, and the predicate body is a single record over canonical JSON~\cite{rfc8785} that binds, in one tuple: treaty id, treaty-scope hash, ladder-intersection hash, admission-report hash, continuation hash, request-canonical hash, outcome-canonical hash, local-receipt hash, remote-receipt hash, capability lease, governance receipt reference, and the two signer kernel ids. The DSSE subject is a single digest whose preimage is the canonical hash of that tuple; the envelope carries exactly two signature slices, one per kernel keyid.

This shape differs structurally from generic DSSE provenance. SLSA provenance and in-toto link statements~\cite{slsa,torres2019intoto} bind a build pipeline to an artifact; their subject is the artifact digest and the predicate records build parameters. The bilateral profile binds a receipt-graph admission event to a treaty, with the subject digest taken over a tuple of hashes rather than an artifact, and with joint authorization intent expressed as two independent kernel signatures over identical canonical bytes. Without that intent record a receiver can verify that some party signed something, but cannot distinguish a shared action from two adjacent local logs that happen to name the same workflow.

\paragraph{The binding tuple.} The subject digest is the canonical hash of a ten-field tuple, and each field commits to a distinct constitutional precondition. Figure~\ref{fig:envelope} visualizes the binding-tuple decomposition. The \emph{treaty-scope hash} pins the bilateral relationship in force (participant identities, named action classes, lifetime), so re-presenting the envelope under a different relationship is a hash miss rather than a re-routable token. The \emph{ladder-intersection hash} pins the joint mode-coverage decision (which action class was evaluated at which mode under the intersected ladder), so a sender cannot retroactively claim a lower-intensity floor than the one both sides agreed to. The \emph{admission-report hash} pins the receiving polity's intent: which predicates the receiver checks, which evidence it requires, and which fail-closed codes it has pre-declared. The \emph{continuation hash} pins the linkage to receipt-graph state by naming the prior receipt this envelope extends, so the envelope cannot be re-grafted onto a different graph position without invalidating the digest. The \emph{request hash} and \emph{outcome hash} commit to the canonical input and output bytes separately, so selective disclosure can later reveal one without the other while preserving the binding. The \emph{local-receipt} and \emph{remote-receipt} hashes carry the sender's and receiver's signed receipt bodies, respectively; the pair is what distinguishes bilateral admission from a unilateral attestation. The \emph{lease} field commits to a freshness window (a revocation epoch and an expiry), and the \emph{signer kernel ids} commit to the two cosigners' identifiers, so an envelope cannot be reassigned to a different pair of signers post-facto.

\begin{figure}[t]
\centering
\begin{tikzpicture}[
  font=\footnotesize,
  envelope/.style={draw, thick, rounded corners=3pt, inner sep=7pt},
  field/.style={draw, rectangle, minimum height=4.2mm, text width=38mm, align=left, font=\scriptsize, inner sep=2pt, fill=blue!4},
  sigbox/.style={draw, rectangle, minimum height=12mm, minimum width=24mm, align=center, font=\scriptsize, inner sep=3pt, fill=green!10},
  hdrlabel/.style={font=\bfseries\scriptsize, align=center},
]
\node[envelope, label={[font=\bfseries\small]above:DSSE Envelope (payload + 2 signatures)}] (env) {%
  \begin{tikzpicture}
    \node[hdrlabel, align=center] (payhead) {payload (canonical JCS bytes)\\\scriptsize\itshape\textmd{bindingTuple = ten fields below}};
    \node[field, below=2mm of payhead.south, anchor=north] (f1) {treaty-scope hash};
    \node[field, below=0.8mm of f1.south, anchor=north] (f2) {ladder-intersection hash};
    \node[field, below=0.8mm of f2.south, anchor=north] (f3) {request digest};
    \node[field, below=0.8mm of f3.south, anchor=north] (f4) {outcome digest};
    \node[field, below=0.8mm of f4.south, anchor=north] (f5) {continuation hash};
    \node[field, below=0.8mm of f5.south, anchor=north] (f6) {local-receipt hash};
    \node[field, below=0.8mm of f6.south, anchor=north] (f7) {remote-receipt hash};
    \node[field, below=0.8mm of f7.south, anchor=north] (f8) {lease (epoch, expiry)};
    \node[field, below=0.8mm of f8.south, anchor=north] (f9) {signer kernel ids ($k_1, k_2$)};
    \node[hdrlabel, right=22mm of payhead] (sighead) {signatures};
    \node[sigbox, below=2mm of sighead.south, anchor=north] (s1) {$\mathrm{sig}_1$\\\scriptsize signed by $k_1$};
    \node[sigbox, below=1.5mm of s1.south, anchor=north] (s2) {$\mathrm{sig}_2$\\\scriptsize signed by $k_2$};
    \node[font=\scriptsize\itshape, align=center, below=2mm of s2.south, text=gray!50!black] {both signatures\\cover the payload bytes};
  \end{tikzpicture}
};
\node[below=2.5mm of env, font=\small] (hash) {$\mathit{subjectDigest} = H(\mathit{canonical}(\mathit{bindingTuple}))$};
\draw[->, thick, shorten >=1pt, shorten <=1pt] (env.south) -- (hash.north);
\end{tikzpicture}
\caption{Bilateral DSSE envelope. The canonical payload binds a ten-field tuple, whose hash under canonical JCS encoding is the subject digest both kernel-key signatures cover.}
\Description{A DSSE envelope diagram with two columns: the left column lists the ten fields of the binding tuple (treaty scope, ladder intersection, request digest, outcome digest, continuation hash, two receipt hashes, lease, signer kernel ids) all contained inside the canonical payload, with a brace grouping them; the right column shows the two signature slots and a note that both signatures cover the payload bytes. Below the envelope, the subject digest equation shows that the subject digest is the hash of the canonical-encoded binding tuple.}
\label{fig:envelope}
\end{figure}

\paragraph{Worked envelope.} A valid envelope, with placeholder hashes, has the canonical shape (each hex field below is the bare lowercase output of SHA-256, truncated here with an ellipsis for display):
\begin{small}
\begin{verbatim}
{ "_type": "in-toto-statement/v1",
  "predicateType":
    "chio.bilateral-cosign-invocation.v1",
  "subject": [ { "digest":
    { "sha256": "9f1c...e0a4" } } ],
  "predicate": {
    "treatyId":         "treaty:opus-alpha:2026",
    "treatyScopeHash":  "7b2a...c1e0",
    "ladderInterHash":  "8c4d...3f02",
    "admissionRptHash": "a91e...204b",
    "continuationHash": "5d22...8e07",
    "requestHash":      "11ff...6a91",
    "outcomeHash":      "33ee...b15d",
    "localReceiptHash": "6677...d4cc",
    "remoteReceiptHash":"8899...02ab",
    "lease": { "epoch": 412, "expiry": 1788000000 },
    "signerKids": [ "kid:opus#1", "kid:alpha#1" ] },
  "signatures": [
    { "keyid": "kid:opus#1",  "sig": "..." },
    { "keyid": "kid:alpha#1", "sig": "..." } ] }
\end{verbatim}
\end{small}
The two \codepath{sig} slices are independent Ed25519 signatures over the identical canonical DSSE pre-authentication encoding of the payload bytes.

\paragraph{Canonical encoding of the binding tuple.} The binding tuple is canonically encoded as an RFC 8785 JCS~\cite{rfc8785} object over a JSON record with ten named fields, one per constitutional precondition: \codepath{treatyScopeHash}, \codepath{ladderInterHash}, \codepath{admissionRptHash}, \codepath{continuationHash}, \codepath{requestHash}, \codepath{outcomeHash}, \codepath{localReceiptHash}, \codepath{remoteReceiptHash}, \codepath{leaseHash}, and \codepath{signerKidsHash}. Field names are fixed strings drawn from this enumeration; each field value is the hex encoding of a SHA-256 output, sixty-four lowercase ASCII characters with no \codepath{sha256:} prefix inside the binding object. The JCS rules force deterministic key ordering, minimal numeric form, and Unicode string normalization, so two implementations that agree on the input fields produce byte-identical encodings. The subject digest is then $\mathrm{SHA\text{-}256}(\mathrm{JCS}(\mathit{bindingObject}))$. The JCS layer supplies field-tag domain separation by construction: two distinct fields cannot collide on attacker-chosen content because their key strings differ in the canonical bytes that feed the digest. Raw concatenation, length-prefixed concatenation with field tags, and HKDF-Expand with field-tag info parameters are not used; the JCS choice inherits the deterministic-serialization guarantee already required for the outer DSSE envelope, and a second canonicalization surface would expose an additional attack surface without strengthening domain separation.

\paragraph{Verifier accept set.} Let the trust store
$\mathcal{S} = (\mathit{IssuerKeys},$
$\mathit{KernelKeys},$
$\mathit{LeaseEpoch})$
and let $E$ be a bilateral envelope. The verifier accepts iff six gates hold in order; (G3) and (G4) below each group two atomic clauses that together express one structural property (signer distinctness, trust-store membership):
\begin{align*}
  \mathit{accept}(\mathcal{S},E) \iff{}
  & \textsf{(G1)}\;\mathit{canonical}(E.\mathit{payload}) \;\wedge \\
  & \textsf{(G2)}\;E.\mathit{predicateType} = \mathit{ChioBilateralV1} \;\wedge \\
  & \textsf{(G3)}\;\bigl(|E.\mathit{sigs}|=2 \;\wedge \\
  & \qquad\;\; E.\mathit{sig}_1.\mathit{kid} \ne E.\mathit{sig}_2.\mathit{kid}\bigr) \;\wedge \\
  & \textsf{(G4)}\;\bigl(E.\mathit{sig}_1.\mathit{kid} \in \mathit{IssuerKeys}\;\wedge \\
  & \qquad\;\; E.\mathit{sig}_2.\mathit{kid} \in \mathit{KernelKeys}\bigr) \;\wedge \\
  & \textsf{(G5)}\;E.\mathit{payload}.\mathit{lease}.\mathit{epoch} \ge \mathcal{S}.\mathit{LeaseEpoch} \;\wedge \\
  & \textsf{(G6)}\;E.\mathit{subjectDigest} = \\
  & \qquad\;\; H(\mathit{canonical}(E.\mathit{payload}.\mathit{bindingTuple})).
\end{align*}
Each verified signature is over the same canonical payload bytes; the verifier never reaches (G1)-(G6) until both bytes-level signatures verify under the keyids declared in the signature slices. The predicate is the conjunction of six gates (G1)-(G6), evaluated left-to-right. Figure~\ref{fig:gates} renders the verifier gate sequence and the five rejection codes. Five of the six gates produce uniquely-named rejection codes; the trust-store-membership gate (G4) is bundled with predicate-type-mismatch in the deployed verifier, because a signer outside both allowlists fails (G2) before (G4) is reached whenever the predicate type is the strict Chio profile and the trust store enumerates only Chio-issuing keys.

\begin{figure}[t]
\centering
\begin{tikzpicture}[
  node distance=3.5mm,
  every node/.style={font=\footnotesize},
  gate/.style={draw, rectangle, rounded corners=2pt, minimum height=6mm, text width=55mm, align=left, inner sep=3pt},
  reject/.style={draw, rectangle, dashed, minimum height=5mm, minimum width=28mm, align=center, font=\scriptsize, fill=red!8, inner sep=2pt},
  accept/.style={draw, rectangle, double, minimum height=6mm, text width=55mm, align=center, fill=green!10},
  arrow/.style={->, >=Stealth, thick},
  rejarrow/.style={->, >=Stealth, dashed, red!70!black},
]
\node[gate] (g1) {(G1) canonical(payload)};
\node[gate, below=of g1] (g2) {(G2) predicateType = ChioBilateralV1};
\node[gate, below=of g2] (g3) {(G3) $|$sigs$|=2 \;\wedge\;$ kids distinct};
\node[gate, below=of g3] (g4) {(G4) kid$_1 \in$ IssuerKeys $\wedge$ kid$_2 \in$ KernelKeys};
\node[gate, below=of g4] (g5) {(G5) lease.epoch $\geq$ S.LeaseEpoch};
\node[gate, below=of g5] (g6) {(G6) subjectDigest = $H$(canonical(bindingTuple))};
\node[accept, below=of g6] (ok) {accept};

\node[reject, right=18mm of g1] (r1) {R1: noncanonical};
\node[reject, right=18mm of g2] (r2) {R2: type mismatch};
\node[reject, right=18mm of g3] (r3) {R3: signer reuse};
\node[reject, right=18mm of g5] (r5) {R4: stale lease};
\node[reject, right=18mm of g6] (r6) {R5: digest mismatch};

\draw[arrow] (g1.south) -- (g2.north);
\draw[arrow] (g2.south) -- (g3.north);
\draw[arrow] (g3.south) -- (g4.north);
\draw[arrow] (g4.south) -- (g5.north);
\draw[arrow] (g5.south) -- (g6.north);
\draw[arrow] (g6.south) -- (ok.north);

\draw[rejarrow] (g1.east) -- (r1.west);
\draw[rejarrow] (g2.east) -- (r2.west);
\draw[rejarrow] (g3.east) -- (r3.west);
\draw[rejarrow] (g5.east) -- (r5.west);
\draw[rejarrow] (g6.east) -- (r6.west);
\end{tikzpicture}
\caption{Verifier gate evaluation. Five named rejection codes distinguish noncanonical payloads, predicate-type mismatch (which bundles G4 trust-store membership), signer reuse, stale leases, and subject-digest mismatch.}
\Description{Vertical flow of six gate boxes labeled G1 through G6 with an accept box at the bottom; dashed arrows from each gate point to red-shaded rejection-code boxes on the right.}
\label{fig:gates}
\end{figure}

\paragraph{noncanonical-payload.} The payload bytes fail RFC 8785 canonicalization: re-ordered keys, non-minimal numeric forms, unnormalized strings, or trailing whitespace. The verifier rejects rather than re-canonicalizing because the signature is over a specific byte sequence. The attack class defeated is the \emph{presentation attacker} who serves different byte sequences to different verifiers, each of which re-canonicalizes locally and reaches a different decision while the upstream signature is undisturbed; the same channel also underwrites equivocation against transparency-log witnesses. Closing it at the verifier protects the canonical-bytes assumption that selective disclosure, replay-detection, and witness consistency all depend on.

\paragraph{predicate-type-mismatch.} The DSSE statement carries a \codepath{predicateType} other than \codepath{chio.bilateral-cosign-invocation.v1}. The attack class defeated is \emph{profile downgrade}: a sender substitutes a weaker predicate (an SLSA build provenance, a generic SCAI record, or a degenerate signed-JSON document) through the same envelope path, hoping the receiver's parser accepts any well-formed DSSE statement and applies bilateral semantics. Generic in-toto profiles do not bind admission-report and ladder-intersection hashes, so admitting them as bilateral collapses joint authorization to unilateral attestation while syntactic envelope checks still pass. Pinning the predicate string at verification rather than at parse time forces the substitution to forge a signature rather than to merely repackage one.

\paragraph{signer-reuse.} The two signature slices carry the same keyid, or distinct keyids that resolve to the same operator under verifier-owned attribution metadata. The attack class defeated is the \emph{confused-deputy cosignature}: one kernel produces both slices, possibly using two of its own key handles, and presents the envelope as bilateral consent. A single signer cannot produce two-party consent; party-independence is what distinguishes bilateral admission from a single kernel signing twice under different aliases. The keyid-equality test catches the trivial form; the operator-resolution test, parameterized by verifier-owned attribution metadata, catches the variant in which one operator runs two kernels under distinct keyids. The gate is conservative: a borderline pair is rejected rather than admitted, since admitting a sole signer destroys the property the envelope is meant to record.

\paragraph{stale-lease.} The capability lease referenced by the payload predates the verifier's current revocation epoch, or its expiry has passed. The attack class defeated is \emph{post-revocation replay}: an attacker holding an envelope whose underlying capability has since been revoked re-presents it to a verifier that has not retained per-action revocation entries, betting that the bilateral signatures remain bytes-valid. The lease binds the action's authority window; admitting a stale lease re-admits a window the issuer has already revoked and turns the receipt graph into an append-anything store. Lease freshness is a verifier-side test against state the verifier owns rather than a sender-asserted claim, which also rules out \emph{epoch downgrade}: the comparison reads the verifier's epoch, not any field in the envelope. Intra-lease replay (resubmitting the same envelope within its still-current lease window) is not prevented at the verifier; it is a constitutional concern handled by the receipt-graph deduplication and continuation-hash discipline the related polity-layer construction develops (anonymized for review), not by the bilateral admission primitive.

\paragraph{subject-digest-mismatch.} The DSSE subject digest is not the canonical hash of the binding tuple constructed from the payload's named fields. The attack class defeated is \emph{sibling-treaty cross-receipt substitution}: an attacker takes a bytes-valid envelope authorized under treaty $\tau_1$ and tries to re-present it under $\tau_2$ by editing the binding tuple's treaty-scope or ladder-intersection field while leaving the upstream signatures in place. Any field-level edit produces a hash miss against the digest the signers committed to. The same gate also rules out \emph{ladder-floor substitution}: the integrity check on the digest is structural, not field-by-field, so any tuple member that disagrees with the committed digest is rejected as a single failure rather than as a per-field diagnostic.

The taxonomy is intentionally coarse and is scoped to the post-signature gates: a signature-bytes failure denies the envelope before the six gates are evaluated and is not a member of the five rejection codes above. Each code names a conjunct of the admission predicate, not a structured failure report identifying which sub-field disagreed at the bit level. An adversary probing the verifier in the non-adaptive setting learns which gate first refused, not which inner field mismatched; the error surface leaks at most $\log_2 5 \approx 2.32$ bits per attempt, and that information is exactly what an auditor needs to identify which constitutional precondition failed. Against an adaptive adversary submitting envelopes in sequence, this bound holds per-attempt but not in aggregate: over $N$ probes against a fixed target envelope, left-to-right gate evaluation reveals only the first refusing conjunct at each step, so an adversary who chains probes can localize the failing gate in $O(N)$ attempts and, with carefully constructed inputs, narrow the disagreement to a particular sub-field. This is the conjunct-localization that a non-error-message-oracle verifier should suppress; doing so requires uniform-time gate evaluation, which the present construction does not provide. Adaptive bounds against the gate-ordering oracle remain open.
